% !TeX root = ../apunte_mecanica_clasica.tex
\IfFileExists{../apunte_mecanica_clasica.tex}{%
  \documentclass[../apunte_mecanica_clasica.tex]{subfiles}%
}{%
  \documentclass[apunte_mecanica_clasica.tex]{subfiles}%
}
\begin{document}

\chapter{Transformaciones de Coordenadas}


\section{Motivación física del cambio de coordenadas}

Cuando estudiamos un sistema físico, anhelamos que las ecuaciones del movimiento sean simples en su forma funcional, con el fin de poder resolverlas.
Pero, a veces no ocurre este anhelo.
Sin embargo, si cambiamos el sistema coordenado a otro, las ecuaciones se podrían, en principio, simplificar.
Por ejemplo, en el problema de fuerzas centrales conservativas, la segunda ley de Newton se reduce a
\begin{equation*}
m\,\vec{a} = f(x^2+y^2)\,\vec{x}.
\end{equation*}
Si se usan coordenadas cartesianas, las ecuaciones del movimiento se escriben de forma explícita como:
\begin{align*}
m\,\ddot{x} &= f(x^2+y^2)\,x, \\[4pt]
m\,\ddot{y} &= f(x^2+y^2)\,y, \\[4pt]
z &= 0.
\end{align*}
Esto se traduce en dos ecuaciones diferenciales ordinarias acopladas y, en general, no lineales. Resolver estas ecuaciones no siempre es simple.

Sin embargo, si usamos coordenadas cilíndricas, 
\begin{equation*} 
\left.\begin{aligned}
x &= r\cos\phi, \\
y &= r\sin\phi, \\
z &= 0, 
\end{aligned} \right\} \qquad r \ge 0,\; 0 \le \phi < 2\pi,
\end{equation*}
las ecuaciones del movimiento del problema central se simplifican a
\begin{align*}
m\!\left(\ddot r - r\,\dot\phi^{\,2}\right) = f(r), \\[4pt]
m\!\left(r^{2}\ddot\phi + 2r\,\dot r\,\dot\phi\right) = 0.
\end{align*}

La segunda ecuación se reduce a
\begin{align*}
\frac{d}{dt}\!\left(m r^{2}\dot\phi\right)=0,
\end{align*}
\begin{equation*}
    \therefore\quad m r^{2}\dot\phi = l.
\end{equation*}

Usando coordenadas cilíndricas se obtiene una constante del movimiento, el momento angular $l$. Por lo tanto, una de las ecuaciones del movimiento reduce su orden diferencial de segundo a primero.

Usando esta constante del movimiento, las ecuaciones del movimiento en coordenadas cilíndricas se reducen a
\begin{align*}
m\,\ddot r \;&=\; f(r) \;+\; \frac{l^{2}}{m\,r^{3}}, 
\\[4pt]
m\,r^{2}\,\dot\phi \;&=\; l.
\end{align*}
Por lo tanto, las ecuaciones del movimiento se simplifican notablemente al pasar de coordenadas cartesianas a cilíndricas. Esto no garantiza su resolubilidad en forma cerrada; sin embargo, la simplificación ayuda de manera sustantiva.


\section{Definición informal y líneas coordenadas}

Una transformación de coordenadas corresponde a un \textbf{cambio de coordenadas que permita escribir las ecuaciones del movimiento en una forma más simple}. Por ejemplo, cambiar de coordenadas cartesianas a coordenadas esféricas
\begin{align*}
x&=r\cos\phi\,\sin\theta, \\[4pt]
y&=r\sin\phi\,\sin\theta, \\[4pt]
z&=r\cos\theta,
\end{align*}
cambio que se traduce en las restricciones
\begin{equation*}
    r\ge 0, \qquad 0\le \phi < 2\pi, \qquad 0\le \theta \le \pi,
\end{equation*}

y en la aplicación
\begin{equation*}
    (x,y,z)\;\longrightarrow\; (r,\phi,\theta).
\end{equation*}

Una transformación de coordenadas no solo cambia los parámetros de las coordenadas, sino que también modifica las líneas coordenadas. En el ejemplo de cambiar de cartesianas a esféricas se cambian de tres rectas a un rayo, una circunferencia y una media circunferencia La figura asociada a esta construcción se muestra en la Figura~\ref{fig:08-transformaciones-de-coordenadas:tdc-fig1}.
\begin{figure}[H]
    \centering
    \includegraphics[width=0.7\linewidth]{figures/16_transformaciones_de_coordenadas/tdc_fig1.png}
    \caption{En el paso de cartesianas a esféricas: con $(\phi,\theta)$ fijos y $r$ variable, la línea coordenada es un rayo que parte del origen; con $(r,\theta)$ fijos y $\phi$ variable, es una circunferencia (paralela al plano $xy$); con $(r,\phi)$ fijos y $\theta$ variable, es una media circunferencia (un meridiano).}

    \label{fig:08-transformaciones-de-coordenadas:tdc-fig1}
\end{figure}


\newpage

\section{Formulación general}

Formalmente, una transformación de coordenadas corresponde a un cambio de la forma
\begin{equation*}
    (x^{1},x^{2},x^{3}) \;\longleftrightarrow\; (x'^{1},x'^{2},x'^{3}),
\end{equation*}

tal que
\begin{equation*}
    \left. \begin{aligned}
        x^{1}=x^{1}(x'^{1},x'^{2},x'^{3})=x^{1}(x'^{j}),\\[6pt]
        x^{2}=x^{2}(x'^{1},x'^{2},x'^{3})=x^{2}(x'^{j}),\\[6pt]
        x^{3}=x^{3}(x'^{1},x'^{2},x'^{3})=x^{3}(x'^{j}),
    \end{aligned}
    \right\} \quad \Rightarrow \quad x^{i}=x^{i}(x'^{j}).
\end{equation*}
El vector posición de un punto arbitrario cambia en la forma:
\begin{equation*}
    \vec{x}(x^{i}) \;=\; \vec{x}\!\big(x^{i}(x'^{j})\big) \;=\; \vec{x}'(x'^{j})
\end{equation*}

\begin{equation*}
    \therefore\quad\vec{x}(x^{i}) \;=\; \vec{x}'(x'^{j}).
\end{equation*}

En general, $\vec{x}(x^{i}) \neq \vec{x}(x'^{\,j})$, puesto que la forma funcional del vector posición puede cambiar con las nuevas coordenadas $x'^{\,i}$. Usaremos la nomenclatura (a menos que se indique explícitamente lo contrario):
\begin{itemize}
    \item $x^{i}$: coordenadas viejas,
    \item $x'^{\,i}$: coordenadas nuevas.
\end{itemize}


\subsection{Ejemplos de cambios de coordenadas}

\begin{apuntebloque}{Ejemplo}

a) Podemos pasar de coordenadas cartesianas a coordenadas cilíndricas
\begin{equation*}
    \left.
    \begin{aligned}
        x^{1}&=x,\\
        x^2&=y,\\
        x^3&=z,
    \end{aligned}
    \right\} \quad \longrightarrow\quad
        \left\{
    \begin{aligned}
        x'^{1}&=r,\\
        x'^2&=\phi,\\
        x'^3&=z.
    \end{aligned}
    \right.
\end{equation*}

a través de la transformación
\begin{align*}
x &= x(r,\phi)=r\cos\phi,\\
y &= y(r,\phi)=r\sin\phi,\\
z &= z(z)=z.
\end{align*}

El vector posición de un punto arbitrario en coordenadas cartesianas y cilíndricas es
\begin{align*}
\vec{x}(x,y,z)&=x\,\hat{\imath}+y\,\hat{\jmath}+z\,\hat{k}\\[4pt]
&= r\cos\phi\,\hat{\imath}+r\sin\phi\,\hat{\jmath}+z\,\hat{k}\\[4pt]
&= \vec{x}'(r,\phi,z),
\end{align*}
donde
\begin{align*}
\vec{x}(x,y,z)&=x\,\hat{\imath}+y\,\hat{\jmath}+z\,\hat{k},\\[4pt]
\vec{x}'(r,\phi,z)&=r\cos\phi\,\hat{\imath}+r\sin\phi\,\hat{\jmath}+z\,\hat{k}.
\end{align*}

En coordenadas cartesianas, el vector posición $\vec{x}(x,y,z)$ es lineal en $x$, $y$ y $z$; en coordenadas cilíndricas, el vector posición $\vec{x}'(r,\phi,z)$ es lineal en $r$ y $z$, y no es lineal en $\phi$ (aparece en funciones trigonométricas).
\\

b) Podemos pasar de coordenadas cartesianas a coordenadas esféricas

\begin{equation*}
    \left.
    \begin{aligned}
        x^{1}&=x,\\
        x^2&=y,\\
        x^3&=z,
    \end{aligned}
    \right\} \quad \longrightarrow\quad
        \left\{
    \begin{aligned}
        x'^{1}&=r,\\
        x'^2&=\phi,\\
        x'^3&=\theta.
    \end{aligned}
    \right.
\end{equation*}

a través de la transformación
\begin{align*}
x &= x(r,\phi,\theta)=r\cos\phi\,\sin\theta,\\
y &= y(r,\phi,\theta)=r\sin\phi\,\sin\theta,\\
z &= z(r,\theta)=r\cos\theta.
\end{align*}

El vector posición de un punto arbitrario en coordenadas cartesianas y esféricas es
\begin{align*}
\vec{x}(x,y,z)&=x\,\hat{\imath}+y\,\hat{\jmath}+z\,\hat{k}\\[4pt]
&= r\cos\phi\sin\theta\,\hat{\imath}+r\sin\phi\sin\theta\,\hat{\jmath}+r\cos\theta\,\hat{k}\\[4pt]
&= \vec{x}'(r,\phi,\theta),
\end{align*}
donde
\begin{align*}
\vec{x}(x,y,z)&=x\,\hat{\imath}+y\,\hat{\jmath}+z\,\hat{k},\\[4pt]
\vec{x}'(r,\phi,\theta)&=r\cos\phi\sin\theta\,\hat{\imath}+r\sin\phi\sin\theta\,\hat{\jmath}+r\cos\theta\,\hat{k}.
\end{align*}

En coordenadas cartesianas, el vector posición $\vec{x}(x,y,z)$ es lineal en $x$, $y$ y $z$; pero, en coordenadas esféricas, el vector posición $\vec{x}'(r,\phi,\theta)$ es lineal solo en $r$ y no es lineal en $\phi$ ni en $\theta$ (aparecen en funciones trigonométricas).

\end{apuntebloque}

\newpage


\section{Coordenizaciones de $\mathbb{R}^3$}

Una coordenización de $\mathbb{R}^3$ es un mapeo de $M\subseteq \mathbb{R}^3$ en $\mathbb{R}^3$, donde
\begin{equation*}
    M = I_1 \times I_2 \times I_3,
\end{equation*}
con $x^{1}\in I_1$, $x^{2}\in I_2$ y $x^{3}\in I_3$, y $\times$ es el símbolo de producto cartesiano. Es claro que $I_i\subseteq \mathbb{R}$, con $i\in \{1,2,3\}$. Luego,
\begin{align*}
\vec{x} &: M \to \mathbb{R}^3, \\[4pt] &\quad x^{i} \mapsto \vec{x}(x^{i}).
\end{align*}

\begin{apuntebloque}{Ejemplo}

a) En coordenadas cartesianas se tiene
\begin{equation*}
    \left.\begin{aligned}
        x^{1}\in I_{1}&=(-\infty,\infty)=\mathbb{R},\\[4pt]
        x^{2}\in I_{2}&=(-\infty,\infty)=\mathbb{R},\\[4pt]
        x^{3}\in I_{3}&=(-\infty,\infty)=\mathbb{R},
    \end{aligned}\right\} \quad
\to \quad M=\mathbb{R}\times \mathbb{R}\times \mathbb{R} = \mathbb{R}^3,
\end{equation*}
por lo tanto
\begin{align*}
\vec{x} &: \mathbb{R}^{3}\to\mathbb{R}^{3}, \\[4pt]  &(x,y,z)\;\mapsto\;\vec{x}(x,y,z).
\end{align*}
\\
b) En coordenadas cilíndricas se tiene
\begin{equation*}
    \left.\begin{aligned}
        r \in I_{1} = [0,\infty),\\[4pt]
        \phi \in I_{2} = [0,2\pi),\\[4pt]
        z \in I_{3} = (-\infty,\infty)=\mathbb{R},
    \end{aligned}\right\} \quad
\to \quad M=[0,\infty)\times[0,2\pi)\times\mathbb{R},
\end{equation*}
por lo tanto
\begin{align*}
\vec{x} &: M \to \mathbb{R}^3, \\[4pt]  
&(r,\phi,z)\;\mapsto\;\vec{x}(r,\phi,z).
\end{align*}
\\

c) En coordenadas esféricas se tiene
\begin{equation*}
    \left.\begin{aligned}
        r \in I_{1} = [0,\infty),\\[4pt]
        \phi \in I_{2} = [0,2\pi),\\[4pt]
        \theta \in I_{3} = [0,\pi],
    \end{aligned}\right\} \quad
\to \quad M=[0,\infty)\times[0,2\pi)\times[0,\pi],
\end{equation*}
por lo tanto
\begin{align*}
\vec{x} &: M \to \mathbb{R}^3, \\[4pt]  
&(r,\phi,\theta)\;\mapsto\;\vec{x}(r,\phi,\theta).
\end{align*}
\end{apuntebloque}


\section{Transformación entre coordenizaciones}

Una transformación de coordenadas es un mapeo de una coordenización en otra de $\mathbb{R}^3$; esto es,
\begin{align*}
x'^{\,i}&: M \to M', \\[4pt]  
&x^{j} \mapsto x'^{\,i}(x^{j}),
\end{align*}
con $x^{i}\in M \;\longleftrightarrow\; x'^{\,i}\in M' \;\longleftrightarrow\; \mathbb{R}^3$.

\newpage

\section{Jacobiano shit}

\subsection{Invertibilidad y matriz jacobiana}

La transformación de coordenadas debe ser invertible, es decir,
\begin{equation*}
    x'^{\,i}(x^{j}) \;\longleftrightarrow\; x^{i}(x'^{\,j}).
\end{equation*}

Por el teorema de la función implícita se sabe que el determinante de la matriz jacobiana de la transformación,
\begin{equation*}
    J^{i}{}_{j} \;=\; \frac{\partial x^{i}}{\partial x'^{\,j}},
\end{equation*}

llamada \textbf{jacobiano} y denotada por
\begin{equation*}
    J \;:=\; \det\!\big(J^{i}{}_{j}\big) \;=\; \det\!\left(\frac{\partial x^{i}}{\partial x'^{\,j}}\right),
\end{equation*}

debe ser distinto de cero,
\begin{equation*}
    J \;\neq\; 0.
\end{equation*}

La matriz jacobiana es
\begin{equation*}
\big(J^{i}{}_{j}\big)\;=\;\left(\frac{\partial x^{i}}{\partial x'^{\,j}}\right)
=\begin{pmatrix}
\dfrac{\partial x^{1}}{\partial x'^{\,1}} & \dfrac{\partial x^{1}}{\partial x'^{\,2}} & \dfrac{\partial x^{1}}{\partial x'^{\,3}}\\[8pt]
\dfrac{\partial x^{2}}{\partial x'^{\,1}} & \dfrac{\partial x^{2}}{\partial x'^{\,2}} & \dfrac{\partial x^{2}}{\partial x'^{\,3}}\\[8pt]
\dfrac{\partial x^{3}}{\partial x'^{\,1}} & \dfrac{\partial x^{3}}{\partial x'^{\,2}} & \dfrac{\partial x^{3}}{\partial x'^{\,3}}
\end{pmatrix}.
\end{equation*}
El índice $i$ rotula las filas y el índice $j$ rotula las columnas.

\begin{apuntebloque}{Definición}
 Una transformación de coordenadas $x'^{\,i}(x^{j})$ se dice \textbf{invertible} si el jacobiano de la transformación es distinto de cero:
\begin{equation*}
J \;:=\; \det\!\big(J^{i}{}_{j}\big)
\;=\; \det\!\left(\frac{\partial x^{i}}{\partial x'^{\,j}}\right)
\;\neq\; 0.
\end{equation*}

\end{apuntebloque}


\subsection{Ejemplos de jacobianos}

\begin{apuntebloque}{Ejemplo}

a) Sea la transformación de coordenadas cartesianas a cilíndricas. Entonces, la jacobiana es
\begin{equation*}
\left(J^{i}{}_{j}\right)
=\left(\frac{\partial x^{i}}{\partial x'^{\,j}}\right)
=\begin{pmatrix}
\dfrac{\partial x}{\partial r} & \dfrac{\partial x}{\partial \phi} & \dfrac{\partial x}{\partial z}\\[8pt]
\dfrac{\partial y}{\partial r} & \dfrac{\partial y}{\partial \phi} & \dfrac{\partial y}{\partial z}\\[8pt]
\dfrac{\partial z}{\partial r} & \dfrac{\partial z}{\partial \phi} & \dfrac{\partial z}{\partial z}
\end{pmatrix}
=\begin{pmatrix}
\cos\phi & -r\sin\phi & 0\\
\sin\phi & \phantom{-}r\cos\phi & 0\\
0 & 0 & 1
\end{pmatrix}.
\end{equation*}

Entonces, el jacobiano de la transformación es
\begin{equation*}
J \;=\; \det\!\big(J^{i}{}_{j}\big)
\;=\; \det\!\left(\frac{\partial x^{i}}{\partial x'^{\,j}}\right)
\;=\; r,
\end{equation*}
el cual es distinto de cero si $r\neq 0$. El caso $r=0$ corresponde al eje $z$ ($x=y=0$, $z\in\mathbb{R}$), donde el vector coordenado $\vec{e}_{\phi}=\partial \vec{x}/\partial \phi$ se anula.

La inversa es
\begin{align*}
r &= \sqrt{x^{2}+y^{2}}\,,\\[6pt]
\phi &= \arctan\!\left(\frac{y}{x}\right)\,,\\[6pt]
z &= z.
\end{align*}

\end{apuntebloque}

\begin{apuntebloque}{}


b) Sea la transformación de coordenadas cartesianas a esféricas. Entonces, la jacobiana es
\begin{equation*}
J^{i}{}_{j}
=\frac{\partial x^{i}}{\partial x'^{\,j}}
=\begin{pmatrix}
\dfrac{\partial x}{\partial r} & \dfrac{\partial x}{\partial \phi} & \dfrac{\partial x}{\partial \theta}\\[8pt]
\dfrac{\partial y}{\partial r} & \dfrac{\partial y}{\partial \phi} & \dfrac{\partial y}{\partial \theta}\\[8pt]
\dfrac{\partial z}{\partial r} & \dfrac{\partial z}{\partial \phi} & \dfrac{\partial z}{\partial \theta}
\end{pmatrix}
=\begin{pmatrix}
\cos\phi\,\sin\theta & -\,r\sin\phi\,\sin\theta & r\cos\phi\,\cos\theta\\
\sin\phi\,\sin\theta & \phantom{-}\,r\cos\phi\,\sin\theta & r\sin\phi\,\cos\theta\\
\cos\theta & 0 & -\,r\sin\theta
\end{pmatrix}.
\end{equation*}

Entonces, el jacobiano de la transformación es
\begin{equation*}
J \;=\; \det\!\big(J^{i}{}_{j}\big)
\;=\; \det\!\left(\frac{\partial x^{i}}{\partial x'^{\,j}}\right)
\;=\; -\,r^{2}\sin\theta,
\end{equation*}
el cual es distinto de cero si $r\neq 0$ y $\theta\neq 0,\pi$. El caso $r=0$ corresponde al origen (un punto), y $\theta=0,\pi$ corresponde al eje $z$ ($x=y=0$, $z\in\mathbb{R}$), donde el vector coordenado $\vec e_{\phi}=\partial \vec x/\partial \phi$ se anula.

La inversa es
\begin{align*}
r &= \sqrt{x^{2}+y^{2}+z^{2}}\,,\\[6pt]
\phi &= \arctan\!\left(\frac{y}{x}\right)\,,\\[6pt]
\theta &= \arccos\!\left(\frac{z}{r}\right)\,.
\end{align*}

\end{apuntebloque}


\newpage

\subsection{Orientación de bases coordenadas}

\textbf{Observación.} La diferencia de signo entre el jacobiano de la transformación a cilíndricas y el de la transformación a esféricas se debe a que las coordenadas cilíndricas poseen vectores coordenados dextógiros, mientras que las coordenadas esféricas poseen vectores coordenados levógiros. La figura asociada a esta construcción se muestra en la Figura~\ref{fig:base_Oxi}.
\begin{figure}[H]
  \centering

  \begin{subfigure}[t]{0.48\textwidth}
    \textbf{Sistema $O\,x^{i}$:}
    \[
      \vec{x}=\vec{x}(x^{i}), 
      \qquad 
      \vec{e}_{i}(x^{j}) \;=\; \frac{\partial \vec{x}(x^{j})}{\partial x^{i}}.
    \]
    \vspace{2mm}
    \centering
    \includegraphics[width=0.8\linewidth]{figures/16_transformaciones_de_coordenadas/tdc_fig2.png}
    \caption{Base coordenada asociada al sistema $O\,x^{i}$.}
    \label{fig:base_Oxi}
  \end{subfigure}
  \hfill
  \begin{subfigure}[t]{0.48\textwidth}
    \textbf{Sistema $O'\,x'^{\,i}$:}
    \[
      \vec{x}'=\vec{x}'(x'^{\,i}), 
      \qquad 
      \vec{e}'_{\,i}(x'^{\,j}) \;=\; \frac{\partial \vec{x}'(x'^{\,j})}{\partial x'^{\,i}}.
    \]
    \vspace{2mm}
    \centering
    \includegraphics[width=1.0\linewidth]{figures/16_transformaciones_de_coordenadas/tdc_fig3.png}
    \caption{Base coordenada asociada al sistema $O'\,x'^{\,i}$.}
    \label{fig:base_Oprime}
  \end{subfigure}

\end{figure}

\newpage

\end{document}
